\documentclass[10pt,a4paper,withhyper,ragged2e]{altacv}
\usepackage{cvstyle}

\begin{document}

\name{Tom Dobson}
\tagline{Data Engineer}
\personalinfo{%
  \email{dobs.tx@gmail.com}
  \location{London, UK}
  \linkedin{dobsontom}
  \github{dobsontom}
  \source{dobsontom/cv}
}
\makecvheader

\cvsection{Summary}

Data engineer owning tails.com's data platform end to end: from the Aurora MySQL and Postgres
application databases to a Snowflake, dbt, Airflow, and Airbyte estate of around 290 billion rows
on AWS and Kubernetes, with the Terraform that provisions it and access managed from group
membership down to per-person database identities. In two years, cut Snowflake spend 46\% and ETL
licence spend 94\% whilst the estate grew, and retired every paid ingestion platform. Inside the
application, profiled the core production database and cut its query load and CPU by around a third
across 2026 through N+1 fixes and indexing. Looking for a senior data engineering role where the
infrastructure is part of the job.

\cvsection{Experience}

\cvrole{Data Engineer}{tails.com}{Mar 2026 -- present}
\begin{itemize}
  \item Cut Snowflake spend \cvmetric{46\%}, saving around \$100k per year, by removing unused and
        over-scheduled compute and sizing each warehouse on measured workload
  \item Migrated every paid ingestion platform to self-hosted Airbyte on EKS, cutting ETL licence
        spend \cvmetric{94\%}, around \$60k per year
  \item Eliminated shared static credentials on production databases, replacing them with
        per-person IAM authentication, 15-minute tokens, and role-scoped grants, so every
        production query traces to an engineer. Live on the core production databases, over
        \cvmetric{4 TB} between them
  \item Designed and built Snowflake governance as code in Terraform, with access managed from
        Google Groups through to Snowflake roles, every service account on key-pair authentication,
        and grants to PUBLIC cut \cvmetric{97\%}
  \item Profiled the production MySQL workload and removed N+1 queries that comprised
        \cvmetric{8\%} of database query volume by caching per-item lookups, and an unindexed join
        scanning around 40,000 rows a second
  \item Built the platform's new ingestion and reverse-ETL pipelines in Airflow, landing survey and
        support-ticket sources in Snowflake and pushing customer segments and orders back into CX
        and marketing tools. Delivered the sync for a new product catalogue API, validated by
        running it in parallel against a clone of production
  \item Protected the analytics estate through a platform migration by auditing every event path
        from the core application into Snowflake, finding the dependency that had no replacement,
        and writing the reference backend engineering used to design the Kafka cart-event topic
        that replaces it
  \item Gave engineers secure AI-agent access to dbt and Snowflake through an MCP server on a
        dedicated identity that can build in each developer's own database but never read unmasked
        PII or write to production
  \item Created and host a biweekly data infrastructure meeting, running since 2025, where analysts
        and developers agree platform changes, audit Snowflake, ingestion, and dbt spend, and
        prioritise the work that follows
\end{itemize}

\cvrole{Analytics Engineer}{tails.com}{Oct 2024 -- Mar 2026}
\begin{itemize}
  \item Owned a dbt project of around \cvmetric{1,600} models on Snowflake and refactored it
        entirely onto dbt's recommended conventions, moving every model into staging, intermediate,
        and marts layers, standardising naming, linting, and grants, and carrying the breaking
        changes through eight repos belonging to other teams
  \item Cut dbt CI time \cvmetric{74\%}, from 14 minutes to around 3, by deferring unchanged models
        to their production tables rather than rebuilding them and capping how far downstream each
        change is tested
  \item Rebuilt deployment and test workflows onto one shared definition of runner, authentication,
        and dbt setup, cutting their YAML \cvmetric{48\%} whilst their number grew, and
        parallelised linting, scoring, and compilation
  \item Brought every model into YAML metadata, up from 21\%, and took documented columns from
        under 1,000 to over \cvmetric{12,000}, with dbt-score failing CI on any touched resource
        that is not fully documented and owned
  \item Removed over \cvmetric{150,000} model builds a month, deprecating 400 dbt models on
        measured build cost and query usage and stopping models rebuilding faster than their data
        changes
  \item Rewrote and incrementalised the most expensive dbt models, cutting the worst from up to 28
        minutes to under \cvmetric{3} by inverting logic that had built every pet and product
        combination, validated row for row against the original
  \item Rebuilt dbt alerting on Elementary, tuned so only the \cvmetric{7.5\%} of detections that
        need action reach an engineer, plus freshness checks that catch silent ingestion failures
  \item Migrated the Airflow estate to v3 and retired v1, and cut image builds \cvmetric{94\%},
        from 7 minutes to under 30 seconds, by stopping the image compiling from source on every
        run
\end{itemize}

\cvrole{BI Consultant, contract via The Information Lab}{Viasat}{Jan 2024 -- Oct 2024}
\begin{itemize}
  \item Lorem ipsum dolor sit amet, consectetur adipiscing elit, sed do eiusmod tempor incididunt
        ut labore et dolore magna aliqua. Ut enim ad minim veniam, quis nostrud exercitation
        ullamco
  \item Duis aute irure dolor in reprehenderit in voluptate velit esse cillum dolore eu fugiat
        nulla pariatur. Excepteur sint occaecat cupidatat non proident
\end{itemize}

\cvrole{Analytics Consultant}{The Information Lab}{Apr 2022 -- Aug 2024}
\begin{itemize}
  \item Sed ut perspiciatis unde omnis iste natus error sit voluptatem accusantium doloremque
        laudantium, totam rem aperiam, eaque ipsa quae ab illo inventore veritatis et quasi
        architecto beatae vitae dicta sunt explicabo. Nemo enim ipsam voluptatem quia voluptas
  \item Ut enim ad minima veniam, quis nostrum exercitationem ullam corporis suscipit laboriosam,
        nisi ut aliquid ex ea commodi consequatur
  \item At vero eos et accusamus et iusto odio dignissimos ducimus qui blanditiis praesentium
        voluptatum deleniti atque corrupti quos dolores et quas molestias
\end{itemize}

\cvsection{Skills}

\cvskillline{Platform}{Lorem ipsum, dolor sit amet, consectetur adipiscing, elit sed do (eiusmod,
  tempor, incididunt), ut labore et dolore, magna aliqua (ut enim, ad minim, veniam, quis),
  nostrud exercitation, ullamco}
\cvskillline{Data and code}{Laboris nisi, ut aliquip (ex ea, commodo, consequat), duis aute irure,
  dolor in, reprehenderit, voluptate velit esse}
\cvskillline{Quality}{Cillum dolore, eu fugiat, nulla pariatur, excepteur, sint occaecat, cupidatat
  non proident, sunt in culpa}
\cvskillline{Prior}{Qui officia, deserunt, mollit anim, id est laborum, sed ut, perspiciatis unde}

\cvsection[nodivider]{Education}

\cvrole{MSc by Research, Biological Sciences}{University of Bristol}{2020 -- 2022}
Thesis nominated for the Postgraduate Dissertation Prize for the best MSc or PhD thesis

\cvrole{BSc (Hons) Biology, First Class}{University of Bristol}{2016 -- 2019}

\cvsection[nodivider]{Training}

\cvrole{AI Bootcamp}{STX Next}{Jun 2026}
Four weeks of hands-on AI-assisted software delivery: planning and context engineering, backend
and frontend delivery, testing, and LLM and MCP integration

\end{document}
